\documentclass[conference]{IEEEtran}
\IEEEoverridecommandlockouts
% The preceding line is only needed to identify funding in the first footnote. If that is unneeded, please comment it out.
\usepackage{cite}
\usepackage{amsmath,amssymb,amsfonts}
\usepackage{algorithmic}
\usepackage{graphicx}
\usepackage{textcomp}
\usepackage{xcolor}
\def\BibTeX{{\rm B\kern-.05em{\sc i\kern-.025em b}\kern-.08em
    T\kern-.1667em\lower.7ex\hbox{E}\kern-.125emX}}
\begin{document}

\title{Conference Paper Title*\\
{\footnotesize \textsuperscript{*}Note: Sub-titles are not captured in Xplore and
should not be used}
\thanks{Identify applicable funding agency here. If none, delete this.}
}

\author{\IEEEauthorblockN{1\textsuperscript{st} Given Name Surname}
\IEEEauthorblockA{\textit{dept. name of organization (of Aff.)} \\
\textit{name of organization (of Aff.)}\\
City, Country \\
email address or ORCID}
\and
\IEEEauthorblockN{2\textsuperscript{nd} Given Name Surname}
\IEEEauthorblockA{\textit{dept. name of organization (of Aff.)} \\
\textit{name of organization (of Aff.)}\\
City, Country \\
email address or ORCID}
\and
\IEEEauthorblockN{3\textsuperscript{rd} Given Name Surname}
\IEEEauthorblockA{\textit{dept. name of organization (of Aff.)} \\
\textit{name of organization (of Aff.)}\\
City, Country \\
email address or ORCID}
\and
\IEEEauthorblockN{4\textsuperscript{th} Given Name Surname}
\IEEEauthorblockA{\textit{dept. name of organization (of Aff.)} \\
\textit{name of organization (of Aff.)}\\
City, Country \\
email address or ORCID}
\and
\IEEEauthorblockN{5\textsuperscript{th} Given Name Surname}
\IEEEauthorblockA{\textit{dept. name of organization (of Aff.)} \\
\textit{name of organization (of Aff.)}\\
City, Country \\
email address or ORCID}
\and
\IEEEauthorblockN{6\textsuperscript{th} Given Name Surname}
\IEEEauthorblockA{\textit{dept. name of organization (of Aff.)} \\
\textit{name of organization (of Aff.)}\\
City, Country \\
email address or ORCID}
}

\maketitle

\begin{abstract}
\end{abstract}

\begin{IEEEkeywords}
\end{IEEEkeywords}

\section{Introduction}

Agent skills are increasingly used to package reusable capabilities for
LLM-based coding agents. In practice, a skill is rarely just a single
prompt file. It often combines instruction text, helper scripts,
configuration, examples, and external resources. This makes it look like
a small software artifact, but one whose interface is partly natural
language and whose operational boundary is entangled with agent runtime
behavior. The central premise of this paper is that skills should be
studied from both a software composition perspective and a security
evaluation perspective. Without the first, security discussions stay too
abstract. Without the second, structural observations remain detached
from practical attack and defense questions.

The paper is organized around three research questions. The first asks
how skills are composed in the wild and what their essential components
are. This question is meant to establish a stable descriptive baseline.
The second asks what attack vectors and attack surfaces are relevant to
skills. This question is intended to refine the current understanding of
skill security by connecting prior work on prompt injection, supply
chain abuse, and agent tooling to the concrete structure of the skill
artifact. The third asks what the current battlefield between attacks
and defenses looks like. This question shifts from description and
taxonomy to comparative evaluation, using representative attacks and
existing scanners to understand which parts of the skill artifact are
well covered and which are still weakly defended.

The full paper will therefore proceed in three layers. It will first
characterize the artifact. It will then organize the threat space around
that artifact. It will finally evaluate how current defensive tools
behave when attacks are instantiated against different skill
components. At the current stage, the purpose of this draft is to fix a
clear and writable outline for those three layers, so that later
iterations can fill in measurements, tables, and discussion without
having to repeatedly change the overall argument.

\section{Background and Scope}

This section will frame skills as a distinct software artifact. The
first part will define what counts as a skill in practice, including the
role of instruction files such as \texttt{SKILL.md}, auxiliary scripts,
metadata, manifests, embedded examples, and references to remote
resources. The goal is to show that a skill is neither purely a prompt
nor purely a package. It is a hybrid artifact whose composition matters
for both usability and security.

The second part of this section will clarify the paper's scope. The
focus will be on attacks that are carried by the skill artifact itself
or that are directly enabled by how the artifact is loaded and used.
This includes malicious instructions, deceptive documentation,
script-level payloads, misleading declarations, and attacks that exploit
the relationship between artifact content and agent action. By contrast,
the paper will not attempt to cover model pretraining attacks, general
web prompt injection unrelated to skills, or broad runtime security
problems that do not depend on the skill channel. This scoping is
important because the space around agent security is already large, and
the paper needs a clear center of gravity.

\section{Methodology}

The methodology section will explain how the three research questions
fit into one coherent study design. The study begins with corpus
construction. We plan to collect roughly one thousand skills from
multiple sources, including public marketplaces and repository-based
distributions, while preserving variation in popularity and publisher
type. The objective is not to build a complete census of the ecosystem
but to construct a corpus that is broad enough to reveal recurring
structural patterns.

After corpus construction, the paper will perform a structural analysis
of the sampled skills. This part supports RQ1. We will inspect file
types, directory layouts, presence of scripts, presence of configuration
artifacts, and references to external dependencies or services. We will
then compare these observations with classic software architecture
concepts such as interface, implementation, dependency, and packaging.
This step is important because the paper is not only asking what skills
contain, but also how that composition should be interpreted.

The next methodological component is a lightweight literature review,
which supports RQ2. The review will use recent work on skill security,
agent security, prompt injection, and transferable supply-chain attacks
as the main search area, followed by one round of backward snowballing
for highly relevant papers and tools. The purpose is not to claim a
complete systematic review. The purpose is to build a defensible and
hierarchical account of attack vectors that is sufficiently grounded in
prior work while remaining directly tied to the skill artifact observed
in RQ1.

The final methodological component is the evaluation protocol for RQ3.
The current benchmark and experimental results will be treated as the
starting point. We will organize existing attack cases by the component
they target, and where necessary we will add manipulated variants that
move the same malicious intent across instructions, scripts, metadata,
and external references. Existing scanners and model-backed analyzers
will then be run under a unified harness. The main output will be a
comparative view of what present defenses can and cannot detect under
the current setup, together with a discussion of how much of that
behavior can be explained by skill composition itself.

\section{RQ1: How Are Skills Composed as a Software Artifact}

The first research question is intended to provide the descriptive
foundation of the paper. This section will begin by quantifying which
major components appear in the sampled corpus. It will examine how often
skills contain only documentation, how often they include executable
files, how often they ship with supporting examples or references, and
how often they declare or imply interaction with external tools and
services. Rather than presenting these observations as an unstructured
inventory, the section will use them to identify a small number of
recurring structural archetypes.

The second part of the section will interpret these archetypes from a
software artifact perspective. A documentation-heavy skill may function
more like a natural-language interface wrapper. A skill with multiple
scripts and auxiliary files may resemble a lightweight plugin package. A
skill with strong dependence on external endpoints may behave more like
a configuration stub that delegates critical behavior elsewhere. The
point of this analysis is to explain what remains familiar from
traditional software packaging and what is genuinely new in the skill
setting.

The final part of RQ1 will set up the later sections by arguing that the
composition of the artifact is not merely descriptive background. It is
also the basis for how attacks can be hidden, how permissions can be
misrepresented, and how scanners may fail. In that sense, RQ1 is the
structural layer that the rest of the paper will depend on.

\section{RQ2: What Are the Attack Vectors and Attack Surfaces for Skills}

The second research question will convert the current flat account of
skill attacks into a more explicit hierarchical structure. This section
will first revisit the most relevant prior work and separate three
sources of attack knowledge: attacks designed specifically for skills,
attacks originally framed for agents or prompt injection but directly
transferable to skills, and attacks adapted from traditional software
supply chains. This separation is necessary because the paper should not
present all attack vectors as equally native to the skill setting.

The next part of the section will organize the attack space by linking
attack goals to artifact components and triggering conditions. At a high
level, the taxonomy will distinguish where the attack is embedded, what
capability it tries to obtain, and under what conditions it becomes
effective. This hierarchical structure is meant to improve over a simple
list of attack names. It should show not only that multiple attack
families exist, but also how they relate to the composition patterns
identified in RQ1.

The final part of RQ2 will explain how this taxonomy will be used later
in the paper. It will serve as the conceptual bridge between corpus
characterization and empirical evaluation. It will define what kinds of
attack placement are meaningful to test, what kinds of scanner coverage
claims can reasonably be made, and where the current literature still
appears incomplete. This makes RQ2 more than a survey section. It is
the organizing theory for the later evaluation.

\section{RQ3: What Does the Current Attack and Defense Battlefield Look Like}

The third research question is the paper's empirical confrontation
section. It will take representative attacks and position them against
current scanners and model-backed analyzers. The purpose is not simply
to produce a leaderboard. The purpose is to observe how defense behavior
changes when attacks are placed in different parts of the skill
artifact, and whether some skill components are consistently easier or
harder to defend than others.

This section will begin by describing the attack set that is already
available and how it will be organized for evaluation. Existing attack
cases provide a practical starting point, but the section will also make
room for manipulated variants that preserve the same malicious intent
while changing the carrier component. For example, an attack originally
expressed in instruction text may be recast through auxiliary scripts or
through misleading metadata. This design will allow the paper to connect
empirical outcomes back to the artifact model developed in RQ1 and the
attack hierarchy developed in RQ2.

The next part will describe the evaluation targets. The plan is to test
multiple scanners and model-backed settings, including less capable
baselines where feasible, under a common protocol. The key question is
not only which system detects more cases overall, but also which kinds
of attacks remain detectable once the same intent is moved across
different components. This is where the paper can make a useful claim
even if absolute scanner performance turns out to be stronger than
initially expected.

The final part of RQ3 will frame the results as a battlefield rather
than a static benchmark report. It will discuss where existing defenses
appear robust, where they appear component-sensitive, and where the
current benchmark remains too weak to support strong claims. This
section should therefore end with a balanced interpretation: either the
evidence supports a meaningful security evaluation paper, or it reveals
that the contribution is narrower and should be framed more modestly.

\section{Discussion}

\section{Conclusion}

\section*{Acknowledgment}

\section*{References}

\begin{thebibliography}{00}
\bibitem{placeholder} References will be inserted after the empirical
sections and literature review are finalized.
\end{thebibliography}

\end{document}
